
\chapter{Resultados da condução da revisão do estado da arte}
\label{ap:resultados}

\section{Relato da condução}

\subsection{Dados gerais da execução}

As buscas foram realizadas em cinco bases de dados científicas
no período de 4 de agosto de 2026 a 11 de outubro de 2026. A
estratégia utilizou termos relacionados ao monitoramento da qualidade
da água, apoio à decisão, dados de sensores e modelos de linguagem
de grande porte.

A string utilizada na maior parte das plataformas foi:

\begin{verbatim}
"water quality monitoring" AND "decision support"
AND "sensor data" AND "large language models"
\end{verbatim}

Não foi utilizado o termo de exclusão \textit{medical}. A data
exata de execução de cada consulta não foi preservada no registro
disponível; portanto, informa-se apenas o período geral conhecido.

A Tabela~\ref{tab:resultados-busca} apresenta a quantidade de registros
informada para cada base de dados.

\begin{table}[htbp]
    \centering
    \caption{Registros retornados pelas bases de dados}
    \label{tab:resultados-busca}
    \begin{tabular}{|l|r|}
        \hline
        \textbf{Base de dados} & \textbf{Registros} \\
        \hline
        Scopus & 15 \\
        \hline
        Web of Science & 22 \\
        \hline
        ScienceDirect & 22 \\
        \hline
        IEEE Xplore & 2 \\
        \hline
        ACM Digital Library & 1 \\
        \hline
        \textbf{Total bruto} & \textbf{62} \\
        \hline
    \end{tabular}
    % -- o\fonte nao ta definido ainda... pensar
    % \fonte{Elaboração própria, com base nos resultados das buscas.}
\end{table}

A soma dos resultados informados pelas cinco bases totalizou
62 registros brutos. Esse total representa a quantidade de resultados
retornados antes da consolidação e da remoção de duplicatas. Assim,
não corresponde necessariamente à quantidade de estudos únicos.

\subsection{Registro das consultas e limitações documentais}

As buscas foram conduzidas ao longo do período informado, mas
não há registro disponível da data exata de cada consulta. Por
esse motivo, não é possível associar uma data individual a cada
base sem introduzir uma informação não documentada.

A string principal foi utilizada na maior parte das plataformas.
As adaptações específicas de sintaxe, os campos pesquisados e
os filtros próprios de cada mecanismo não estão integralmente
documentados nesta versão. Essa limitação deverá ser considerada
na avaliação da reprodutibilidade do processo de busca.

Os endereços das plataformas foram utilizados para consultar
os resultados. Entretanto, os endereços disponíveis para IEEE
Xplore e ACM Digital Library correspondem às páginas gerais
das plataformas, não constituindo, isoladamente, evidência
persistente da consulta específica realizada.

\subsection{Processo de seleção dos estudos primários}

Os 62 registros identificados constituem o conjunto inicial de
resultados da revisão. A partir desse conjunto, devem ser realizadas
a consolidação dos registros, a identificação de duplicatas e a
avaliação da elegibilidade das publicações.

Nesta versão, as quantidades de duplicatas, exclusões por título
e resumo, textos completos avaliados e estudos incluídos ainda
não foram consolidadas. Esses valores não são apresentados como
resultados porque dependem da execução e do registro das etapas
de seleção.

A Figura~\ref{fig:fluxo-revisao} representa, de forma provisória,
a etapa inicial documentada do processo.

\begin{figure}[htbp]
    \centering
    \fbox{
        \parbox{0.85\textwidth}{
            \centering
            Identificação nas cinco bases de dados

            \vspace{0.3cm}

            Scopus: 15 registros

            Web of Science: 22 registros

            ScienceDirect: 22 registros

            IEEE Xplore: 2 registros

            ACM Digital Library: 1 registro

            \vspace{0.3cm}

            Total inicial: 62 registros

            \vspace{0.3cm}

            Próximas etapas a concluir: remoção de duplicatas,
            triagem por título e resumo, avaliação dos textos
            completos e definição do conjunto final de estudos.
        }
    }
    \caption{Etapa inicial documentada do processo de seleção}
    \label{fig:fluxo-revisao}
    % -- o\fonte nao ta definido ainda... pensar
    %\fonte{Elaboração própria.}
\end{figure}

\section{Tabela dos estudos primários selecionados}

A tabela definitiva dos estudos primários deverá ser preenchida
após a conclusão da seleção. Como essa etapa ainda não foi
consolidada, não são apresentados títulos, autores ou anos como
se constituíssem o conjunto final de estudos incluídos.

Após a seleção, a tabela deverá conter, no mínimo, o título
da publicação, os autores e o ano de publicação. Recomenda-se
também registrar a base de origem e um identificador persistente,
quando disponível.

\section{Análise descritiva}

A análise descritiva será realizada sobre o conjunto final de
estudos selecionados. Como a seleção ainda não foi concluída,
não são apresentados gráficos de distribuição temporal, veículos
de publicação, países, instituições ou autores, pois tais análises
exigem dados bibliográficos consolidados.

\subsection{Distribuição dos artigos no tempo}

Após a seleção dos estudos primários, será elaborado um gráfico
com a quantidade de publicações por ano. A análise buscará
identificar a evolução temporal das pesquisas relacionadas ao
tema da revisão.

\subsection{Tipos de veículos de publicação}

Os estudos selecionados serão classificados de acordo com
o tipo de veículo de publicação, como periódicos científicos
e conferências. A distribuição será apresentada em gráfico
e acompanhada de uma interpretação dos resultados.

\subsection{Distribuição por país, instituição e autores}

Serão analisadas as informações bibliográficas relativas aos
países, às instituições e aos autores dos estudos selecionados.
Os resultados serão apresentados por meio de gráficos ou tabelas
adequados à distribuição observada, com o objetivo de identificar
concentrações de publicações e padrões de colaboração.